\documentclass[11pt]{article}
\usepackage[letterpaper,margin=1in]{geometry}
\usepackage{amsmath,amssymb}
\usepackage{booktabs}
\usepackage{graphicx}
\usepackage{caption}
\usepackage{xcolor}
\usepackage[colorlinks=true,linkcolor=blue,citecolor=blue,urlcolor=blue]{hyperref}
\captionsetup{font=small,labelfont=bf}

\newcommand{\muC}{\mu_{\mathcal{C}}}
\newcommand{\muQ}{\mu_{\mathcal{Q}}}
\newcommand{\capture}[1]{A_k(#1;\,\mathcal{C},d)}
\newcommand{\todo}[1]{\textcolor{red}{[TODO: #1]}}

\title{\textbf{The Bond Metric: Breaking Apparent Trade-offs in\\
High-Dimensional Embedding Geometry}\\[4pt]
\large A measure--metric--scale decomposition, a falsification chain,\\
and why one-number hubness is non-identifying}
\author{Andrew H. Bond,~\emph{Senior Member, IEEE}\\
Department of Computer Engineering, San Jos\'e State University\\
\texttt{andrew.bond@sjsu.edu} ~\textbar~ ORCID: 0009-0003-2599-6158}
\date{Working draft --- v0.1 (skeleton; Round-7 verdict pending)}

\begin{document}
\maketitle

\begin{abstract}
Summary statistics of embedding geometry --- intrinsic dimension, effective rank,
hubness skew, ball-growth slope --- are routinely treated as properties of a corpus
and are widely believed to trade off against one another. We show that several of
these trade-offs are \emph{representational artifacts}: they appear when a
generative mechanism (or a learned representation) binds distinct geometric layers
to a single control, and they dissolve when the layers are given separate
mechanisms. We propose a four-layer decomposition of embedding geometry --- global
metric/spectral structure, corpus measure $\muC$, query measure $\muQ$, and local
scale structure --- and support it with a pre-registered falsification chain over a
procedural corpus-generator search (seven experiment rounds against a real
$1024$-dimensional text-embedding corpus). Three constructive separations result:
(i) angular hub structure vs.\ global spectrum, an apparently sharp Pareto
frontier dissolved by splitting placement from colouring; (ii) \emph{same hubness,
different anatomy}: a synthetic corpus and the real corpus share the same scalar
hubness reading produced by opposite mechanisms (centrality-driven super-hubs,
max reverse-$k$NN count 369, vs.\ density-driven moderate hubs, max count 78) ---
one-number hubness is non-identifying; (iii) hubness vs.\ ball growth, dissolved
by recognising that the real system's hubness lives substantially in the query
marginal: writing hubness as $H_k(\mathcal{C},\mathcal{Q},d)$, reverse-neighbour
counts sample the query mass of top-$k$ capture basins,
$N_k(x_j)\approx|\mathcal{Q}|\,\muQ(\capture{x_j})$. We conclude that apparent
geometric trade-offs can be artifacts of forcing corpus metric, corpus measure,
query measure, and local scale through a shared mechanism, propose a hub
\emph{anatomy vector} as the reporting standard, and derive a registered
anatomy falsifier that prevents optimizers from reproducing a scalar by
rebuilding a known pathology. \todo{Round-7 verdict sentence.}
\end{abstract}

\section{Introduction}\label{sec:intro}
% Narrative arc: six rounds spent solving a corpus-side problem that lived
% substantially in the query marginal.
\begin{itemize}
\item The setting: procedural generators that must reproduce real embedding
  geometry (benchmark admission gates G1--G8 of OpenVector Bench; RC-1/RC-2
  discipline), searched adversarially with pre-registered predictions and named
  falsifiers --- every round's misses reported as misses.
\item The mystery: successive families each solved gates individually
  (intrinsic dimension, rank, PCA retention, local-ID spread, hubness) yet hit
  hard pairwise ``walls'': hubness $\times$ rank (round 5), then hubness
  $\times$ ball growth (round 6--7 diagnosis).
\item The resolution: the walls were properties of \emph{mechanism binding},
  not of the statistics. Contributions:
  \textbf{(C1)} the four-layer decomposition (\S\ref{sec:decomposition});
  \textbf{(C2)} three constructive separations with pre-registered protocols
  (\S\ref{sec:separations});
  \textbf{(C3)} the capture-basin account of hubness,
  $H_k(\mathcal{C},\mathcal{Q},d)$, and the non-identifiability of scalar
  hubness (\S\ref{sec:hubness});
  \textbf{(C4)} the anatomy vector $A_{\mathrm{hub}}$ and the registered
  anatomy falsifier (\S\ref{sec:anatomy});
  \textbf{(C5)} implications for ANN benchmark design and differential system
  testing (\S\ref{sec:implications}; companion paper: \emph{Fuzzing the Bond
  Metric}).
\item Epistemic status, stated plainly: the decomposition is \emph{strongly
  corroborated and materially sharpened} by the falsification chain --- not
  universally confirmed; a general independence theorem would require formal
  construction under stated assumptions (\S\ref{sec:limits}).
\end{itemize}

\section{The four-layer decomposition}\label{sec:decomposition}

\begin{table}[h]\centering\small
\begin{tabular}{lll}
\toprule
Layer & Main observable & Control \\
\midrule
Global metric / spectral & spectrum, effective rank, PCA retention & recolouring \\
Corpus measure $\muC$ & cluster mass, density, angular capture & centre placement \\
Query measure $\muQ$ & visitation frequency of corpus regions & query marginal \\
Local scale structure & ball growth, neighbourhood transitions & radius/scale spectrum \\
\bottomrule
\end{tabular}
\caption{The Bond decomposition. ``Measure'' is not one thing: two interacting
measures, $\muC$ and $\muQ$, act on the same metric substrate. The query measure
is an \emph{illumination pattern} over the corpus's capture basins.}
\label{tab:layers}
\end{table}

\begin{itemize}
\item Formal statement of each layer; which registered gates (G1--G8) read
  which layer, and which read mixtures. \todo{gate-to-layer map table}
\item The binding thesis: a trade-off observed under one generative family is
  evidence about that family's control structure, not about the statistics.
\item Amplification: in local dimension $d_{\ell}$, measure-layer observables
  respond as $s^{-d_{\ell}}$ to metric-layer perturbations of relative scale
  $s$ --- binding layers is easy to do by accident and expensive to undo.
\end{itemize}

\section{Experimental substrate and protocol}\label{sec:substrate}
\begin{itemize}
\item OpenVector Bench: the eight geometry gates, battery A (corpus$\to$corpus)
  vs.\ battery B (held-out queries$\to$corpus); real target = BGE-M3--class
  $1024$-d text embeddings, $n=8$k--$16$k slices; sealed 25\% never touched.
\item The search harness: \texttt{evaluate\_fn} contract, 12-shard swarm random
  search (480 evals/round) on NRP Nautilus; searcher $\times$ adversary
  $\times$ judge; pre-registration with named falsifiers before every gating
  run; screening smokes disclosed in the registration when they occurred.
\item Round-by-round provenance table (rounds 0--7, one line each: family,
  registered prediction, verdict). \todo{table from results/ directory}
\end{itemize}

\section{Three constructive separations}\label{sec:separations}

\subsection{Separation I: angular hubs vs.\ global spectrum (rounds 5--6)}
\begin{itemize}
\item Round 5: Zipf codebook hierarchy controls both hub placement and the
  covariance spectrum $\Rightarrow$ clean two-ended Pareto frontier: every
  hub-matching shard has effective rank $\le 0.05\times$; the one rank-matching
  shard sits at hubness $0.24\times$; nothing in the middle. Registered
  falsifier fired exactly as written.
\item Round 6: same placement mechanism \emph{inside} a Mahalanobis
  recolouring of the centred covariance (separate control) $\Rightarrow$ the
  frontier dissolves; first five-gate joint match
  (G1/G3/G6/G7/G8 all within $[0.5,2.0]\times$, stable across 3 seeds).
  Recolouring \emph{amplifies} rather than destroys the hubs.
\item Table: round-5 endpoints vs.\ round-6 joint point. \todo{numbers from
  gen\_round5\_hier.json / gen\_round6\_coloured.json}
\end{itemize}

\subsection{Separation II: same hubness, different anatomy (the headline)}
\begin{itemize}
\item Hub-anatomy measurement, real corpus vs.\ the round-6 winner at matched
  scalar G6: real --- base$\to$base skew $\approx 1.5$, max reverse-count 78,
  density-driven ($\mathrm{corr}(N_{10}, -d_{10}) = +0.67$), hub $d_{10}$ only
  8\% below population; generated --- skew $7.0$, max count \textbf{369},
  centrality-driven, hubs 34\% closer than population.
\item One-number hubness is non-identifying: two systems share G6 with opposite
  geometry and (predictably) opposite ANN behaviour. \todo{ANN behaviour
  demonstration --- candidate for the Fuzzing paper instead}
\item Table: the anatomy vector for real 8k/16k, Gaussian null, round-6 winner.
  \todo{numbers from diag\_hubs.json}
\end{itemize}

\subsection{Separation III: query measure vs.\ local scale (round 7)}
\begin{itemize}
\item The G2 mechanics: the ball-growth estimator's radius window is the
  across-query $d_{10}$ band; slope $\approx \log 9 / \log(\mathrm{window})$;
  real window $1.17 \Rightarrow$ slope $\approx 15$. \todo{derivation box}
\item Six corpus-side hub mechanisms fail on a metrically homogeneous substrate
  (equalization, mean offset, duplicate families, per-point tilt, in-patch
  sub-clusters, cluster count) --- reported as the negative results they are.
\item The diagnosis: real base$\to$base skew $\approx1.5$ vs.\ battery-B target
  $6.8$ --- the hubness lives in $\muQ$. A same-instance Zipf query marginal
  sweeps G6 $0.19\times \to 2.20\times$ with the window intact and G2
  \emph{improving} to $0.65\times$.
\item The old G6--G2 conflict, unwound: same-distribution queries cannot
  produce query-induced G6 $\to$ search grows corpus-side central super-hubs
  $\to$ extreme density contrast $\to$ patch/plateau anatomy $\to$ G2
  collapses. The conflict was between \emph{the wrong mechanism for hubness}
  and realistic ball growth.
\item \todo{ROUND-7 VERDICT: six-gate joint point (P1), anatomy falsifier
  (P2), G8 risk (P3) --- fill from GEN\_ROUND7\_QUERY.md}
\end{itemize}

\section{Hubness is $H_k(\mathcal{C},\mathcal{Q},d)$: the capture-basin account}
\label{sec:hubness}
\begin{itemize}
\item Definition: $\capture{x_j}$ = the region of query space for which $x_j$
  ranks in the top $k$ under metric $d$ against corpus $\mathcal{C}$. Then
  \begin{equation}
  N_k(x_j) \;\approx\; |\mathcal{Q}| \cdot \muQ\!\left(\capture{x_j}\right).
  \label{eq:capture}
  \end{equation}
\item Corpus geometry defines the basins; the query marginal decides how
  heavily each is visited; identical skew can arise from huge basins
  (corpus-side pathology) or concentrated query mass (the real mechanism).
\item Consequences for the hubness literature: base$\to$base hubness
  (Radovanovi\'c-style, $H_k(\mathcal{C},d)$) and query-measured hubness are
  different observables and should not be conflated; benchmark suites that
  generate queries from the corpus distribution structurally cannot exhibit
  the query-induced regime. \todo{position against related work}
\end{itemize}

\section{The anatomy vector and the registered anatomy falsifier}\label{sec:anatomy}
\begin{itemize}
\item Reporting standard:
  \[
  A_{\mathrm{hub}} = (\text{q$\to$b skew},\ \text{b$\to$b skew},\ \max N_k,\
  \mathrm{corr}(N_k, d_M),\ \mathrm{corr}(N_k, -d_{10}),\ \text{top-hub mass share}).
  \]
\item The registered defense (round-7 P2): the optimizer's best point must keep
  base$\to$base skew $<3$ --- the first pre-registered falsifier that looks
  \emph{behind} a gate at the structure producing it, preventing scalar
  reproduction by rebuilding a known pathology. \todo{verdict}
\item Generalization: every scalar gate deserves an anatomy check; sketch the
  analogous vectors for G1 (near-duplicate spike vs.\ manifold bulk) and G3
  (spectrum shape vs.\ participation ratio).
\end{itemize}

\section{Implications}\label{sec:implications}
\begin{itemize}
\item For benchmark design: corpus-only generators cannot reproduce
  query-induced statistics; workload models need an explicit $\muQ$.
\item For ANN evaluation: same-G6 corpora with opposite anatomy are natural
  differential tests (exact vs.\ HNSW/DiskANN/FAISS/cuVS; fp32 vs.\ quantized;
  metric choice; build order; sharding), with oracles on hub-rank divergence
  and anti-hub-query recall collapse --- aggregate recall hides exactly the
  tail the decomposition exposes. Companion paper: \emph{Fuzzing the Bond
  Metric} (coverage signature over layer-separated cells).
\item For representation learning: apparent statistic trade-offs in learned
  embeddings deserve the same suspicion --- they may be facts about the
  architecture's control structure, not about geometry.
\end{itemize}

\section{Related work}\label{sec:related}
\begin{itemize}
\item Hubness: emergence with dimensionality, centrality accounts
  (Radovanovi\'c et al.), hubness reduction (local scaling, mutual proximity)
  --- reread as metric-layer vs.\ measure-layer interventions. \todo{cites}
\item Query/corpus distribution shift in IR and ANN benchmarks. \todo{cites}
\item Intrinsic-dimension estimation (TwoNN, Levina--Bickel) and its
  sensitivity to near-duplicates and sampling. \todo{cites}
\item Coverage-guided fuzzing and property-based testing (the methodological
  transfer). \todo{cites}
\end{itemize}

\section{Limitations and the road to a theorem}\label{sec:limits}
\begin{itemize}
\item Corroboration, not universal confirmation: three separations on one real
  corpus family (text embeddings, one model class) and one generator lineage.
\item What a theorem needs: a formal generative class in which the four layers
  are independently parameterized, plus proofs that each gate is invariant to
  the other layers' controls within stated ranges; candidate assumptions
  (unit-sphere metric, sub-Gaussian patches, capture-basin regularity).
\item Known couplings that are real, not representational: the $s^{-d_\ell}$
  amplification; normalization's coupling of radial to angular structure;
  where the decomposition should be expected to fail.
\end{itemize}

\section{Conclusion}\label{sec:conclusion}
\todo{Write after round-7 verdict. Closing line candidate: six rounds were
spent solving a corpus-side problem whose target lived substantially in the
query marginal --- the statistics were never in conflict; the mechanisms were.}

\paragraph{Reproducibility.} All rounds, pre-registrations, falsifier verdicts,
and raw numbers: \url{https://github.com/ahb-sjsu/openvector-bench}
(\texttt{results/}, \texttt{spec/BOND\_METRIC.md}). Sealed evaluation rows were
never loaded during any experiment reported here.

\end{document}
